\section{The \texorpdfstring{$N=1$}{N=1} Oberdieck and Pixton normalisation}
\label{wn:sec:lorgat-conj1-N1-OP-normalization}

\providecommand{\phiZero}{\phi_{0,1}}
\providecommand{\gDelta}{\mathfrak{g}_{\Delta_5}}
\providecommand{\ZDTred}{Z^{\mathrm{red}}_{\mathrm{DT}}}
\providecommand{\Dfive}{D_5}
\providecommand{\PhiOP}{\Phi_{10}^{\mathrm{OP}}}

Let \(X=K3\times E\).  Let
\(\beta_h\in H_2(K3,\mathbb Z)\) be primitive effective with
\(\langle\beta_h,\beta_h\rangle=2h-2\).  Write the Eichler and Zagier
normalised Jacobi form as in \emph{The Theory of Jacobi Forms},
Table~1:
\[
  \phiZero(\tau,z)=\sum_{a,b} f(a,b)q^a y^b
  =
  (y^{-1}+10+y)
  +
  q(10y^{-2}-64y^{-1}+108-64y+10y^2)+O(q^2).
\]
The full K3 elliptic genus is \(Z_{K3}=2\phiZero\).  Thus \(f(0,0)=10\),
while the full K3 elliptic genus has zeroth coefficient \(20\).

Let \(\alpha=(n,\ell,m)\in\Lambda^{2,1}_{\mathrm{II}}\) be a positive
root of \(\gDelta\), put \(h=nm\), and assume \(\gcd(n,\ell,m)=1\).

\begin{theorem}[Primitive exponent and scalar square]
\label{thm:lorgat-conj1-N1-OP-normalization}
In the primitive Gritsenko and Nikulin chamber,
\[
  \operatorname{sdim}\mathfrak g_{\Delta_5,\alpha}=f(nm,\ell).
\]
The Oberdieck and Pixton scalar partition function is the reciprocal
square
\[
  \ZDTred(q,\tilde q,y)
  =
  -(\PhiOP)^{-1}
  =
  -\Dfive^{-2}
  =
  -4096\,\Delta_5^{-2},
  \qquad
  \Dfive=64^{-1}\Delta_5,\quad
  \PhiOP=\Dfive^2=4096^{-1}\Delta_5^2.
\]
Equivalently,
\[
  -\ZDTred
  =
  q^{-1}r^{-1}s^{-1}
  \prod_{(a,b,c)>0}(1-q^ar^b s^c)^{-2f(ac,b)} .
\]
The Borcherds modular characteristic is
\[
  \kappa_{\mathrm{BKM}}(\Delta_5)
  =
  \left.c_N(0)/2\right|_{N=1}
  =
  10/2
  =
  5 .
\]
\end{theorem}

\begin{proof}
Borcherds' singular theta lift gives weight \(c(0)/2\) for the
automorphic product attached to a weight zero Jacobi input
\cite[Theorem~13.3]{Borcherds1998}.  For \(\phiZero\), the zeroth
coefficient is \(f(0,0)=10\), so the output weight is \(5\).
Gritsenko and Nikulin identify the \(N=1\) product with
\[
  \Dfive(Z)=64^{-1}\Delta_5(Z)
  =
  q^{1/2}r^{1/2}s^{1/2}
  \prod_{(a,b,c)>0}(1-q^ar^b s^c)^{f(ac,b)}
\]
\cite[Theorem~2.1]{GritsenkoNikulin1998}.  The denominator identity of
the BKM superalgebra reads the exponent of the primitive root
\(\alpha=(n,\ell,m)\) as the signed root superdimension
\(\operatorname{sdim}\mathfrak g_{\Delta_5,\alpha}=f(nm,\ell)\).

Oberdieck and Pixton use the full K3 elliptic genus in the Igusa
product; its exponents are \(2f(nm,\ell)\).  Therefore their monic
Igusa product is
\[
  \PhiOP=\Dfive^2=(64^{-1}\Delta_5)^2=4096^{-1}\Delta_5^2.
\]
Oberdieck and Pixton, \emph{Inventiones Mathematicae}~213 (2018),
Theorem~1, prove the primitive reduced series identity
\[
  Z^X_{\mathrm{OP}}=-(\PhiOP)^{-1},
\]
in the chamber \(0<|\tilde q|\ll |q|\ll |y|^{-1}\ll1\).  Oberdieck,
\emph{Mathematische Zeitschrift}~289 (2018), Theorem~3(i), identifies
the reduced DT and PT series for non-zero K3 class.  After the usual OP
variable convention, this scalar series is \(\ZDTred\).  The displayed
scalar identity follows.

The primitive BKM statement concerns logarithmic Euler product
exponents.  A coefficient of \(\ZDTred\) is the plethystic convolution
of all reciprocal-square factors of lower and equal degree.  Thus the
protected one-particle datum is \(f(nm,\ell)\) on the square-root side
and \(2f(nm,\ell)\) on the OP scalar side; scalar DT coefficients are
read from the reciprocal product.
\end{proof}

\paragraph{Catalogue separation.}
The primitive CHL denominator ladder has
\[
  (c_N(0))_{N=1,2,3,4,6}=(10,8,6,4,2),
  \qquad
  \bigl(\kappa_{\mathrm{BKM}}(\Phi_N)\bigr)=(5,4,3,2,1).
\]
This is the BKM denominator scope of Gritsenko~1999, Theorem~1.2,
together with Borcherds' weight formula.
The Gritsenko and Clery diagonal divisor catalogue is a separate
triple-indexed list with twice-weight tuple
\[
  (10,4,6,2,4,1,3,2).
\]
Its classification and weight formula are Gritsenko and Clery~2008,
Theorems~1.4 and~3.1.
The common row is \(F_{1,1}=\Delta_5=\Phi_1\).  The present note uses
only this common row and the OP scalar square of its monic product.

\paragraph{\texorpdfstring{\(K3\times E\)}{K3 x E} invariants.}
For the compact threefold \(X=K3\times E\),
\[
  \kappa_{\mathrm{cat}}(X)=0,\qquad
  \kappa_{\mathrm{ch}}(X)=0,\qquad
  \kappa_{\mathrm{ch}}^{\mathrm{Heis}}(X)=3,\qquad
  \kappa_{\mathrm{BKM}}(\Delta_5)=5,\qquad
  \kappa_{\mathrm{fiber}}(K3)=24.
\]
These are invariants of distinct constructions: Hodge supertrace,
Heisenberg specialisation, Borcherds weight, and Mukai fibre rank.

\paragraph{Imprimitive roots.}
For \(\gcd(n,\ell,m)>1\), the divisor-sum normalisation used for the
imprimitive BKM character is
\[
  \operatorname{sdim}\mathfrak g_{\Delta_5,\alpha}
  =
  \sum_{d\mid\gcd(n,\ell,m)}
  \mu(d)\,f(nm/d^2,\ell/d).
\]
The reduced stable-pair extension in imprimitive K3 class is governed
by the corresponding multiple-cover formula.  The theorem above is the
primitive \(N=1\) square-root and OP-normalisation statement.
